%% Chapter 6 — Section 6.1: Unadjusted Langevin Algorithm
\section{Unadjusted Langevin Algorithm}
\label{sec:6.1}

(ULA) relies on a stochastic variation of the deterministic update \eqref{eq:6.1} to obtain
approximate samples from a target distribution $f(x) \propto \exp(-V(x))$, see Figure~\ref{fig:6.1}.
Starting from a random variable $X^{(0)}$, draws $\{X^{(n)}\}_{n=1}^{\infty}$ from ULA satisfy
\begin{equation}
  \label{eq:6.1-update}
  \mathbb{E}\bigl[X^{(n+1)} \mid X^{(n)}\bigr] = X^{(n)} - \epsilon \nabla V\!\bigl(X^{(n)}\bigr),
  \qquad n \ge 0.
\end{equation}
Therefore, on average, ULA moves $X^{(n)} \mapsto X^{(n+1)}$ follow the direction of
$-\nabla V(X^{(n)})$, decreasing the value of $V$ and increasing the value of the target
$f \propto \exp(-V)$ provided that the step-size is sufficiently small.
The method is summarized in Algorithm~\ref{alg:6.1}.

\begin{figure}[htbp]
  \centering
  \includegraphics[width=0.75\textwidth]{figures/fig_06_1}
  \caption{\textit{Objective function $V(x) = x^4 - x^2 - 0.4x$ and corresponding target density
    $f(x) \propto \exp(-V(x))$. Low values of $V$ correspond to large values of $f$.}}
  \label{fig:6.1}
\end{figure}

\begin{algorithm}[H]
  \caption{Unadjusted Langevin Algorithm (ULA)}
  \label{alg:6.1}
  \begin{algorithmic}[1]
    \Require Target $f(x) \propto \exp(-V(x))$, initial distribution $\pi_0$,
             step-size $\epsilon > 0$, sample size $N$.
    \State \textbf{Initial draw:} Sample $X^{(0)} \sim \pi_0$.
    \State \textbf{Subsequent draws:} For $n = 0, 1, \ldots, N-1$ set
    \[
      X^{(n+1)} = X^{(n)} - \epsilon\nabla V\!\bigl(X^{(n)}\bigr) + \sqrt{2\epsilon}\,\xi^{(n)},
      \qquad \xi^{(n)} \sim \Normal(0, I).
    \]
    \Ensure Sample $\{X^{(n)}\}_{n=1}^{N}$.
  \end{algorithmic}
\end{algorithm}

\begin{remark}[Pros and Cons]
ULA can be implemented without knowing the normalizing constant of the target distribution.
Moreover, by taking steps $X^{(n)} \mapsto X^{(n+1)}$ that are, on average, in the direction
of $-\nabla V(X^{(n)})$, ULA is able to quickly find and explore regions of high target density.
Under convexity assumptions on $V$ and for sufficiently small step-size $\epsilon$, draws from
ULA are, after a burn-in period, approximately distributed according to the target.
However, the algorithm is biased in that, for fixed $\epsilon$, the distribution of the $n$-th
sample does not converge to the target in the large $n$ asymptotic.
Another caveat of ULA is that taking steps in the direction of $-\nabla V(X^{(n)})$ can be
counterproductive if the target is multi-modal. In that case, ULA can get stuck in one mode
without adequately exploring others. Finally, ULA requires evaluating $\nabla V$, which can
be expensive.
\end{remark}

We will analyze ULA for Gaussian target $f = \Normal(\mu, \Sigma)$. In this setting, similar
to the Gibbs sampler in Chapter~5, draws from ULA satisfy a simple recursion. As in Chapter~5,
we will denote by $H = \Sigma^{-1}$ the precision matrix.

\begin{lemma}
\label{lem:6.1}
Let $\{X^{(n)}\}_{n \ge 1}$ be the output of ULA with target $f = \Normal(\mu, \Sigma) =
\Normal(\mu, H^{-1})$ and initial state $X^{(0)}$. It holds that
\begin{equation}
  \label{eq:6.2}
  X^{(n+1)} - \mu = (I - \epsilon H)^{n+1}(X^{(0)} - \mu)
  + \sqrt{2\epsilon}\sum_{i=0}^{n}(I - \epsilon H)^{n-i}\xi^{(i)}.
\end{equation}
Suppose further that $\pi_0 = \Normal(\mu_0, \sigma_0^2 I)$ for some $\mu_0 \in \R^d$ and
$\sigma_0^2 > 0$. Then, $X^{(n)} \sim \Normal(\mu_n, \Sigma_n)$, where
\begin{align}
  \mu_n &= \mu + (I - \epsilon H)^n(\mu_0 - \mu), \label{eq:6.3} \\
  \Sigma_n &= \Bigl(H - \tfrac{\epsilon}{2}H^2\Bigr)^{-1}
    + (I - \epsilon H)^{2n}\Bigl(\sigma_0^2 I - \Bigl(H - \tfrac{\epsilon}{2}H^2\Bigr)^{-1}\Bigr).
    \label{eq:6.4}
\end{align}
\end{lemma}

\begin{proof}
Due to the Gaussian target assumption, we can set $V(x) = \frac{1}{2}(x-\mu)^T H(x-\mu)$,
which implies that $\nabla V(x) = H(x-\mu)$. Therefore,
\begin{align*}
  X^{(n+1)} - \mu
  &= X^{(n)} - \epsilon\nabla V\!\bigl(X^{(n)}\bigr) + \sqrt{2\epsilon}\,\xi^{(n)} - \mu \\
  &= (I - \epsilon H)(X^{(n)} - \mu) + \sqrt{2\epsilon}\,\xi^{(n)} \\
  &= (I - \epsilon H)^{n+1}(X^{(0)} - \mu) + \sqrt{2\epsilon}\sum_{i=0}^{n}(I-\epsilon H)^{n-i}\xi^{(i)},
\end{align*}
thus establishing \eqref{eq:6.2}. Under the additional assumption that
$\pi_0 = \Normal(\mu_0, \sigma_0^2 I)$, it follows from \eqref{eq:6.2} that each $X^{(n)}$
has Gaussian distribution. To compute the mean, using that $\mathbb{E}[X^{(0)}] = \mu_0$ and
that $\mathbb{E}[\xi^{(i)}] = 0$ for all $i \ge 0$,
\begin{align*}
  \mu_n - \mu
  &= \mathbb{E}\Bigl[(I-\epsilon H)^n(X^{(0)}-\mu)
    + \sqrt{2\epsilon}\sum_{i=0}^{n-1}(I-\epsilon H)^{n-i-1}\xi^{(i)}\Bigr] \\
  &= (I - \epsilon H)^n(\mu_0 - \mu),
\end{align*}
which proves \eqref{eq:6.3}. To compute the covariance, note first that
\begin{align*}
  X^{(n)} - \mu_n
  &= (I-\epsilon H)^n(X^{(0)}-\mu) + \sqrt{2\epsilon}\sum_{i=0}^{n-1}(I-\epsilon H)^{n-i-1}\xi^{(i)}
    - (\mu_n - \mu) \\
  &= (I-\epsilon H)^n(X^{(0)}-\mu_0) + \sqrt{2\epsilon}\sum_{i=0}^{n-1}(I-\epsilon H)^{n-i-1}\xi^{(i)}.
\end{align*}
Hence, using that $\mathbb{E}[(X^{(0)}-\mu_0)(X^{(0)}-\mu_0)^T] = \sigma_0^2 I$, that
$\mathbb{E}[\xi^{(i)}(\xi^{(i)})^T] = I$, and the independence of $X^{(0)}$ and each
$\xi^{(i)}$ from all other randomness, we deduce that
\begin{align*}
  \Sigma_n
  &= \mathbb{E}\bigl[(X^{(n)}-\mu_n)(X^{(n)}-\mu_n)^T\bigr] \\
  &= \sigma_0^2(I-\epsilon H)^{2n}
    + 2\epsilon\sum_{i=0}^{n-1}(I-\epsilon H)^{2i} \\
  &= \Bigl(H - \tfrac{\epsilon}{2}H^2\Bigr)^{-1}
    + (I-\epsilon H)^{2n}\Bigl(\sigma_0^2 I - \Bigl(H-\tfrac{\epsilon}{2}H^2\Bigr)^{-1}\Bigr),
\end{align*}
where in the last line we used the formula for the partial sum of the geometric series
$\sum_{i=0}^{n-1} A^i = (I-A)^{-1} - A^n(I-A)^{-1}$. We have hence established \eqref{eq:6.4},
completing the proof.
\end{proof}

The key takeaway from Lemma~\ref{lem:6.1} is that, for fixed and sufficiently small $\epsilon$,
we have
\begin{align*}
  \mu_n &\xrightarrow{n\to\infty} \mu, \\
  \Sigma_n &\xrightarrow{n\to\infty} \Sigma_\epsilon
    := \Bigl(H - \tfrac{\epsilon}{2}H^2\Bigr)^{-1}.
\end{align*}
The next theorem shows that, in Wasserstein distance, $\pi_n := \Normal(\mu_n, \Sigma_n)$
converges exponentially to $f_\epsilon := \Normal(\mu, \Sigma_\epsilon)$ and that $f_\epsilon$
is order $\epsilon$ apart from $f$. We recall that the Wasserstein distance between Gaussians
$\pi = \Normal(\mu,\Sigma)$ and $\tilde\pi = \Normal(\tilde\mu, \tilde\Sigma)$ is given by
\[
  W_2(\pi, \tilde\pi)^2 = \norm{\mu - \tilde\mu}^2 + \norm{\Sigma^{1/2} - \tilde\Sigma^{1/2}}_F^2
\]
provided that $\Sigma$ and $\tilde\Sigma$ commute. Here $\norm{A}_F = \sqrt{\operatorname{Tr}(A^T A)}$
denotes the Frobenius norm. We will denote by $\lambda_{\min}(H)$ and $\lambda_{\max}(H)$ the
smallest and largest eigenvalues of $H$.

\begin{theorem}[Convergence of ULA for Gaussian Targets]
\label{thm:6.2}
Let $\{X^{(n)}\}_{n\ge 1}$ be the output of ULA with target $f = \Normal(\mu,\Sigma) =
\Normal(\mu, H^{-1})$ and initial distribution $\pi_0 = \Normal(\mu_0, \sigma_0^2 I)$.
Let $\pi_n$ denote the distribution of $X^{(n)}$.
Suppose that $\sigma_0^2 \le \lambda_{\max}(H)^{-1}$.
Then, there is a constant $c$ which depends on $\mu_0, \Sigma_0, \mu$, and $\Sigma$ such that
\[
  W_2(\pi_n, f) \le c\bigl(1 - \epsilon\lambda_{\min}(H)\bigr)^n
  + \frac{\epsilon}{4}\sqrt{\operatorname{Tr}(H)} + \mathcal{O}(\epsilon^2),
  \qquad n = 1, 2, \ldots
\]
\end{theorem}

\begin{proof}
By triangle inequality, $W_2(\pi_n, f) \le W_2(\pi_n, f_\epsilon) + W_2(f_\epsilon, f)$.
The first term represents the distance between the distribution of $X^{(n)}$ and the limit
distribution of ULA, and the second term represents the distance between the limit distribution
of ULA and the target. We will bound both terms in turn. We denote by $c$ a constant that
depends on $\mu_0, \Sigma_0, \mu$, and $\Sigma$ and which may change from line to line.

To bound $W_2(\pi_n, f_\epsilon)$, we need to control the deviation between the means and
covariances of $\pi_n = \Normal(\mu_n, \Sigma_n)$ and $f_\epsilon = \Normal(\mu, \Sigma_\epsilon)$.
For the means,
\begin{align*}
  \norm{\mu_n - \mu}^2
  &= \norm{(I-\epsilon H)^n(\mu_0 - \mu)}^2 \\
  &\le \max\bigl\{|1-\epsilon\lambda_{\min}(H)|,\,|1-\epsilon\lambda_{\max}(H)|\bigr\}^{2n}
     \norm{\mu_0-\mu}^2 \\
  &= c\bigl(1 - \epsilon\lambda_{\min}(H)\bigr)^{2n},
\end{align*}
where we used the Cauchy--Schwarz inequality and that, for all $\epsilon$ sufficiently small,
$\max\{|1-\epsilon\lambda_{\min}(H)|,|1-\epsilon\lambda_{\max}(H)|\} = |1-\epsilon\lambda_{\min}(H)|$.

For the covariances, let $Q\operatorname{diag}(h_1,\ldots,h_d)Q^T$ be the eigendecomposition
of $H$, where $\{h_i\}_{i=1}^d$ are the eigenvalues of $H$ and $Q$ is orthonormal. Then,
\[
  \Sigma_\epsilon = \Bigl(H - \tfrac{\epsilon}{2}H^2\Bigr)^{-1}
    = Q\operatorname{diag}(\alpha_1,\ldots,\alpha_d)Q^T, \quad
  \Sigma_n = Q\operatorname{diag}(\alpha_1-\beta_1,\ldots,\alpha_d-\beta_d)Q^T,
\]
where $\alpha_i = \bigl(h_i - \frac{\epsilon}{2}h_i^2\bigr)^{-1}$ and
$\beta_i = (1-\epsilon h_i)^{2n}(\alpha_i - \sigma_0^2)$.
Notice that, for all sufficiently small $\epsilon$, we have that $\frac{\epsilon}{2}h_i < 1$,
and hence that $\alpha_i > 0$.
Furthermore, the assumption on $\sigma_0^2$ ensures that $\sigma_0^2 \le \alpha_i$,
and hence that $\beta_i \ge 0$.
Consequently, for all $\epsilon$ sufficiently small,
\begin{align*}
  \norm{\Sigma_n^{1/2} - \Sigma_\epsilon^{1/2}}_F^2
  &= \sum_{i=1}^d \bigl((\alpha_i-\beta_i)^{1/2} - \alpha_i^{1/2}\bigr)^2
   = \sum_{i=1}^d \bigl(\alpha_i - \beta_i + \alpha_i - 2\sqrt{\alpha_i(\alpha_i-\beta_i)}\bigr) \\
  &\le \sum_{i=1}^d \bigl(\alpha_i - \beta_i + \alpha_i - 2\sqrt{(\alpha_i-\beta_i)(\alpha_i-\beta_i)}\bigr)
   = \sum_{i=1}^d \beta_i \\
  &= \sum_{i=1}^d (1-\epsilon h_i)^{2n}(\alpha_i - \sigma_0^2)
   \le c\bigl(1-\epsilon\lambda_{\min}(H)\bigr)^{2n},
\end{align*}
where for the last inequality we use that $\operatorname{Tr}(\Sigma_\epsilon - \Sigma_0) \to
\operatorname{Tr}(\Sigma - \Sigma_0)$ as $\epsilon \to 0$, which implies that
$\operatorname{Tr}(\Sigma_\epsilon - \Sigma_0)$ is uniformly bounded for all sufficiently small
$\epsilon$. Combining the bounds for the means and the covariances,
\begin{equation}
  \label{eq:6.5}
  W_2(\pi_n, f_\epsilon)^2 \le c\bigl(1 - \epsilon\lambda_{\min}(H)\bigr)^{2n}.
\end{equation}

Next, we bound $W_2(f_\epsilon, f)$. Now the means agree, and we deduce as before that
\begin{align*}
  W_2(f_\epsilon, f)^2
  &= \sum_{i=1}^d \frac{1}{h_i}\Bigl(1 - \frac{1}{\sqrt{1-\frac{\epsilon}{2}h_i}}\Bigr)^2
   = \sum_{i=1}^d \frac{1}{h_i}\,\varphi\!\Bigl(\tfrac{\epsilon}{2}h_i\Bigr),
\end{align*}
where $\varphi(z) := \bigl(1 - \frac{1}{\sqrt{1-z}}\bigr)^2$.
By Taylor expansion around zero, one can show that, for small $z$,
$\varphi(z) = \frac{z^2}{4} + \mathcal{O}(z^3)$. Hence, for small $\epsilon$,
\begin{equation}
  \label{eq:6.6}
  W_2(f_\epsilon, f)^2
  \le \frac{\epsilon^2}{16}\sum_{i=1}^d h_i + \mathcal{O}(\epsilon^3)
  = \frac{\epsilon^2}{16}\operatorname{Tr}(H) + \mathcal{O}(\epsilon^3).
\end{equation}
Taking square roots in \eqref{eq:6.5} and \eqref{eq:6.6} and using triangle inequality
completes the proof.
\end{proof}

\paragraph{Choice of Step-Size}
The proof of Theorem~\ref{thm:6.2} decomposes the Wasserstein distance between $\pi_n$ and
$f$ into two terms. The first term, which corresponds to the distance between $\pi_n$ and
$f_\epsilon$, decreases exponentially fast with the number $n$ of iterations. The second term,
which corresponds to the distance between $f_\epsilon$ and $f$, is a systematic bias of order
$\epsilon$. Notice that a smaller step-size reduces the bias, but it makes slower the convergence
of $\pi_n$ to the biased limit $f_\epsilon$. Thus, it is recommended to take the largest
$\epsilon$ such that the error in the approximation $f \approx f_\epsilon$ can be tolerated.
While not considered here, time-varying step-sizes can also be used; in particular, a common
approach is to take bold large steps for a few iterations to ensure fast exploration, and then
gradually decrease the step-size to reduce the bias.
